\subsection{Costo total de la Implementación de Mejora / Innovación / Creatividad}

Se presentara un cuadro donde se detallara el costo total del desarrollo del software.

\par
\noindent
\begin{minipage}{\linewidth}
    \tituloTablaAPA{tab:costos_totales_proyecto}{Resumen general de costos totales de implementación del proyecto}

    \begingroup
    \renewcommand{\baselinestretch}{1}
    \footnotesize
    \setlength{\tabcolsep}{4pt}
    \renewcommand{\arraystretch}{1.25}

    \begin{tabularx}{\linewidth}{@{}p{6.5cm}r r@{}}
        \toprule
        \textbf{Descripción} & \textbf{Total (S/)} & \textbf{Dólares (USD)} \\ \midrule
        Costo de Materiales & S/ 876.00 & \$254.65 \\ 
        Costo de Mano de Obra & S/ 71,400.00 & \$20,755.81 \\ 
        Costo Maquinaria, herramientas y equipos & S/ 2,253.66 & \$655.13 \\ 
        Otros Costos (Servicios básicos e Internet) & S/ 869.34 & \$252.72 \\ \midrule
        \textbf{Total} & \textbf{S/ 75,399.00} & \textbf{\$21,918.31} \\ \bottomrule
    \end{tabularx}
    \endgroup

    \notaTabla{Resumen consolidado del costo total de inversión y desarrollo del proyecto. El rubro de maquinaria representa la depreciación anual de los 6 equipos de cómputo. Valores calculados al tipo de cambio referencial de S/ 3.44 por dólar (fecha: 09/11/2026). Elaboración propia.}
\end{minipage}
\finBloqueAPA

Se concluyew que la implementación del sistema tiene una inversión total de S/ 75,399.00 (USD 21,918.31) considerando los costos de mano de obra, materiales, lincecias de software, que son necesarias para el desarrollo del proyecto.